\section{Related Work}\label{sec:related}
We group prior studies into classical diffusion and latent-space models, quantum generative and denoising models, and work on the design, trainability and benchmarking of quantum circuits. For each study we state the problem, method, findings and limitations, and how it relates to this work.

Ho et al.~\cite{ho} addressed high-quality image synthesis with denoising diffusion probabilistic models, which learn to reverse a gradual noising process using a weighted variational bound. They reported a FID of 3.17 on unconditional CIFAR-10. The approach relies on many sequential denoising steps and is entirely classical. It provides the denoising objective that motivates our work, although we study only a single step at one noise level. To reduce the high training and sampling cost of pixel-space diffusion, Rombach et al.~\cite{rombach} ran the diffusion process in the latent space of a pretrained autoencoder and added cross-attention for conditioning. Their models required far less computation than pixel-based ones while remaining highly competitive across several synthesis tasks. Generation is still iterative, and quantum components are not considered. We adopt the principle of denoising a compact latent representation. As a convolutional baseline we use a lightweight network based on the U-Net~\cite{unet}, an encoder--decoder with skip connections originally proposed for biomedical image segmentation.

K\"olle et al.~\cite{kolle} investigated whether variational circuits can reduce the parameter demands of diffusion models. They proposed a densely entangled circuit, a quantum U-Net and a single-step sampling scheme. At about 1,000 parameters on $8\times8$ MNIST, the dense circuit outperformed classical dense networks of equal size, although a much larger U-Net performed best overall. The authors note limits on fidelity due to the normalisation of measured outputs, slow simulation and reliance on noise-free training. Their equal-parameter comparison motivates the matched controls that our study still lacks. Cacioppo et al.~\cite{cacioppo} used parameterized circuits as the denoising operator of a diffusion model, including a variant operating on autoencoder latent vectors. On MNIST, the latent variant achieved far better FID than the fully quantum one, and a reduced version ran on IBM hardware. The full quantum model was limited to $16\times16$ images of two digits, and the comparison with classical models was mainly qualitative. This work is closely related to ours in combining an autoencoder with a quantum denoiser; we instead process overlapping spatial patches and report classical baselines on the same test data.

Zhang et al.~\cite{zhang} studied generative learning of ensembles of quantum states. Their quantum denoising diffusion probabilistic model scrambles states in the forward process and learns the reverse process step by step with shallow circuits, which outperformed their quantum GAN and direct-transport baselines on a clustered-state task. Experiments were simulated with only a few qubits and concerned quantum rather than classical data. The work shows the value of dividing denoising into smaller steps, which our single-step model does not yet exploit. Parigi et al.~\cite{parigi} asked whether quantum noise can benefit generative diffusion, and proposed three combinations of classical or quantum forward and reverse processes. Simulations on two-dimensional points and single-qubit states showed accurate reconstruction. Scalability and any quantum advantage remain to be demonstrated. Our setting corresponds to their classical-noise, quantum-denoising case, applied to image latents. Huang et al.~\cite{huang} addressed image generation on near-term hardware with a flexible quantum GAN. They generated handwritten digit images on a superconducting processor and reported performance competitive with classical GANs on a grey-scale bar dataset. The benchmark is simple and the model is adversarial rather than a denoiser. Its relevance lies in handling high-dimensional images with limited quantum resources, which our patch decomposition also targets.

Turning to circuit design, Havl\'{\i}\v{c}ek et al.~\cite{havlicek} used quantum states as feature spaces for classification, implementing a variational classifier and a kernel estimator on a superconducting processor, and reported successful classification of synthetic data. Their study concerns classification, and its usefulness for image denoising is not addressed. It informs our IQP-style encoding. P\'erez-Salinas et al.~\cite{salinas} showed that repeatedly re-uploading data between trainable operations lets even a single qubit act as a universal classifier when combined with a classical subroutine, and that entangled qubits improve efficiency. Schuld et al.~\cite{schuld} explained such effects theoretically: a variational quantum model can be written as a partial Fourier series whose accessible frequencies are set by the encoding gates, so that repeated encoding enriches the representable functions. Neither study isolates the effect of re-uploading in a hybrid image denoiser. These results motivate the six-qubit design of V2.

McClean et al.~\cite{mcclean} examined the trainability of randomly initialised circuits and showed that, for a broad class of them, the probability of a non-negligible gradient decreases exponentially with the number of qubits. The result does not directly predict the behaviour of small structured hybrid models, but it cautions against assuming that larger circuits will train well. Bowles et al.~\cite{bowles} benchmarked 12 quantum models on 160 datasets. Off-the-shelf classical models generally performed better, and removing entanglement often gave equal or better results. The study is restricted to small binary classification tasks, many of them artificial. We apply its emphasis on strong baselines and component removal to latent denoising.

Table~\ref{tab:related} summarises the most closely related studies. Together they show that quantum circuits can be trained for generative and denoising tasks and that encoding and benchmarking choices are decisive. Their settings differ from ours in data, task or evaluation protocol. PHQD examines patch-wise latent denoising, compares several circuit designs against classical and linear baselines on one test set, and measures the effect of removing the circuit output. A parameter-matched classical substitute remains future work.

\begin{table}[t]
\caption{Summary of closely related studies.}\label{tab:related}
\centering\scriptsize
\renewcommand{\arraystretch}{1.25}
\begin{tabular}{@{}>{\raggedright\arraybackslash}p{0.20\textwidth}>{\raggedright\arraybackslash}p{0.23\textwidth}>{\raggedright\arraybackslash}p{0.26\textwidth}>{\raggedright\arraybackslash}p{0.23\textwidth}@{}}
\toprule
Study & Method & Main finding & Limitation for our problem\\
\midrule
Rombach et al.~\cite{rombach} & Classical diffusion in autoencoder latent space & Much lower computation with competitive quality & No quantum component\\
K\"olle et al.~\cite{kolle} & Variational circuits in diffusion models & Better than equal-size dense networks on small MNIST & Larger U-Net still best; noise-free training\\
Cacioppo et al.~\cite{cacioppo} & PQC denoiser, including latent variant & Latent variant much better than fully quantum & Small scale; limited classical comparison\\
Zhang et al.~\cite{zhang} & Step-wise quantum diffusion & Outperforms quantum baselines on a state-learning task & Quantum data; few qubits\\
Parigi et al.~\cite{parigi} & Classical or quantum noise and denoising & Accurate reconstruction on toy data & Toy scale; no advantage shown\\
Huang et al.~\cite{huang} & Quantum GAN on hardware & Generates digit images & Simple benchmark; not a denoiser\\
P\'erez-Salinas et al.~\cite{salinas} & Data re-uploading & Single qubit acts as universal classifier & Classification only\\
Bowles et al.~\cite{bowles} & Large-scale benchmarking & Classical models generally stronger & Classification only\\
\bottomrule
\end{tabular}
\end{table}
